Learning Outcomes
Upon completion of the course, students will be able to select adequate computational techniques and apply appropriate computational models and algorithms to complex biological problems and assess the range of validity of each.

Course Contents/ Curriculum
Computational methods for solving biological problems and designing synthetic life-like systems at various scales. 
Methods including physics-based approaches such as particle-based atomistic and mesoscopic simulations, as well as techniques in data-driven bioinformatics and machine learning. 
Physics-based approaches include recent advances in Monte Carlo, molecular dynamics, and Brownian dynamics simulations, as well as kinetic modeling.
Data-driven techniques for analyzing next generation sequencing experiments, including transcriptome and single cell analysis. 
Topics are guided by examples from current research advances and challenges from recent literature or faculty research. For each case study topic dealing with a specific subset of computational techniques, the relevant physical, chemical, or mathematical principles are discussed. Explanatory material on the case study, relevant background, and a code or software example will be distributed prior to class. Practical applications in a computer laboratory complement the lectures. In the practical part, depending on the complexity of the computer-based method, (pseudo) code examples are developed in class or supplemented with critical components. Scientific software is also used in practical exercises to solve the case study problem. The range and possible pitfalls of the applied methods are critically examined.

Indications Assignments:

 Form teams of two students each. Indicate the matriculation number of the team members, before July 8 at 17:00, here.  
Select one of the following assignments (description will be attached soon).
Assignment 1: The compaction of intrinsically disordered proteins studied by machine learning
Assignment 2: MD simulation of a protein with different force fields
Assignment 3: Correlation KO and compound inhibition in HTS experiments
Assignment 4: Source publicly available data to interpret HTS results   
Select the project of your preference in the poll (link will be available on July 8 at 18:00)
First come first served basis
In the poll under "Your name" please write 
Matriculation number 1 - Matriculation number 2   (NO STUDENT NAMES)

ONLY ONE VOTE PER TEAM

Upon consolidation of the results from the poll, We will distribute the full assignments by July 10.
Submission of the assignment:  one assignment per team should be submitted by July 26 (single pdf file, 2-6 pages). 